\caption{\textbf{Regional distribution of response types for different numbers of k-means clusters.} k-means as in Fig.~2 (k = 25 reproduces the paper's labels exactly) for k = 10--100; 158 Beryl regions with at least 50 neurons. Black, real regions; red, regions of the same sizes formed by shuffling neurons between regions (100 shuffles), which gives what sampling alone allows. Lines, median over regions; bands, interquartile range (b: 95\% range over shuffles). \textbf{a} Number of response types present in a region, as a fraction of k. \textbf{b} Fraction of region $\times$ type cells whose share of the region lies within two-fold of the type's brain-wide share. \textbf{c} Specialization $S = 1 - H(p)/\log k$ (Fig.~3). \textbf{d} Spearman correlation across regions between $S$ at each k and $S$ at k = 25.
\label{fig:k_robustness}}
